\clearpage
\fancyhead[L]{\small LINDENAU. ZACH.}
\hypertarget{native-490-lower}{}
\label{native-490-lower}
\begin{center}
\rule{0.55\linewidth}{0.4pt}\\[2pt]
Monatliche Correspondenz. Band XIV. Seite 56. 57. 1806 Juli.\\[2pt]
\rule{0.55\linewidth}{0.4pt}
\end{center}

\begin{center}
Ueber \textsc{La Place}'s Ausdruck für Höhenmessungen durch Barometer. Aus einem Schreiben eines \emph{alten} Recensenten [v. \textsc{Lindenau}.]
\end{center}

\begin{flushright}
--- --- --- 1806 Mai 18.
\end{flushright}

--- --- Ueber Ungereimtheiten fällt mir \textsc{Rohde}'s letzte in der M.\,C. recensirte Abhandlung wieder bei, über welche ich und N.\,N. in unsern Urtheilen differirten und wo N.\,N. die ganze Methode auch theoretisch für falsch erklärte, was ich nicht that. Hierüber schreibt er [\textsc{Gauss}] mir nun folgendes:

`Auf einen in Ihrem vorletzten Briefe geäusserten Zweifel muss ich noch einmal zurückkommen. Sie finden nemlich meine Gründe, warum ich \textsc{Rohde}'s Methode auch theoretisch für falsch erklärt habe, zwar einleuchtend, indessen äussern Sie doch, dass der Umstand, dass \textsc{Rohde}'s Ausdruck mit einem bekannten übereinstimme, Sie bedenklich gemacht hätte, die Methode unbedingt für falsch zu erklären.
\SourceLeaf{491}
Mir deucht, dass hier blos ein kleiner Missverstand zum Grunde liege. Ich habe nemlich nicht die von \textsc{Rohde} gebrauchte Formel für falsch erklärt, sondern sein Unternehmen mit Hülfe einer zwar wahren Formel ein Element zu bestimmen. was sich durch diese Formel nicht bestimmen lässt.'

Hiemit bin ich nun auch mit N. N. einverstanden. Ich hätte diesen Umstand nicht erwähnt, wenn Hauptmann \textsc{Rohde} nicht wieder Gelegenheit dazu gegeben hätte etc.

\clearpage
\hypertarget{native-491-lindenau}{}
\label{native-491-lindenau}
\begin{center}
[VON LINDENAU AN GAUSS.]

[In \textsc{Gauss} Nachlass.]
\end{center}

\begin{flushright}Altenburg den 17. Febr. 1806.\end{flushright}

--- --- --- Ich theilte den Brief, in dem Ew. Wohlgeboren mir Ihr Gutachten über die \textsc{Rohden}'sche Antikritik zu communiciren die Güte hatten, Hrn. Obh. \textsc{von Zach} mit, der sich, wie Sie aus der M. C. werden ersehen haben, veranlasst fand, den ganzen Brief einzurücken, was, wie ich hoffe, Ihrer Willensmeinung nicht entgegen sein wird. --- ---

Noch muss ich eines kleinen Zweifels erwähnen, der mir bei dieser Gelegenheit beigegangen ist. Ew. Wohlgeb. erklären \textsc{Rohdens} Methode auch in \emph{theoretischer} Hinsicht für falsch, und die von Ihnen hiefür beigebrachten Gründe sind mir allerdings ganz einleuchtend. Bedenke ich dagegen auf der andern Seite, dass \textsc{Rohdens} Ausdruck, wie ich in meiner Beantwortung gezeigt habe, mit dem bekannten Verhältniss zwischen Massen, Umlaufszeiten und mittlern Entfernungen, analog ist, so machte mich dies wieder bedenklich jene Methode unbedingt für falsch zu erklären. --- ---

\clearpage
\hypertarget{native-491-zach}{}
\label{native-491-zach}
\begin{center}
VON ZACH AN GAUSS.
\end{center}

\begin{flushright}Eisenberg den 6. April 1807.\end{flushright}

Sie haben mich, verehrungswürdigster Freund, durch die gefällige Uebersendung der \textsc{Güssmann}schen Broschüre unendlich aber noch mehr durch Ihre gründliche Beurtheilung verpflichtet. Ich werde solche in das unter der Presse befindliche Maiheft, jedoch wie Sie wünschen, mit Verschweigung Ihres Namens bringen, so gern ich meinerseits gewünscht hätte diesen ausstellen zu dürfen; allein natürlich muss und werde ich jederzeit Ihren Willen und Ihre Ruhe respectiren. Ich bin schon zufrieden, dass Sie meiner Bitte Gehör gegeben, und meine Vorstellung gegründet finden, dass man so böse und boshafte Menschen, wie die Socii unseres P. \textsc{Güssmann}'s sind, in ihrer wahren Gestalt darstellen muss, um manchen verdienten Mann bei Ehren zu erhalten, und bei nicht gehörig unterrichteten Personen vor Nachtheil und Schaden zu bewahren. Ein solches verdienstliches und gutes Werk haben Sie, Verehrtester, nicht nur für die gute Sache, sondern auch für Ihre guten Freunde unternommen. Nicht nur diese, sondern alle wohldenkenden Menschen werden Ihnen dieses ewig Dank wissen. --- --- ---
